\chapter{Probability Quantifiers: Uniform Finite-Model Semantics}
\label{sec:p15-probability-quantifiers}

This is the finite-model presentation of \(L_{\omega P}\) in
Keisler--Lotfallah's author-hosted manuscript, especially its model
conventions on p.~3 and opening probability-quantifier treatment on
pp.~5--7 \cite{KeislerLotfallah2009AuthorProbability}. It is a distinct
presentation, not a reconstruction of Keisler's complete 1985 chapter
on probability quantifiers or its conditional-expectation operators.
The following semantic and translation proofs are supplied editorially.

\begin{hbdefinition}{def:p15-pq-models}{Two measures with different sample spaces}
Fix a nonempty finite relational vocabulary \(\nu\), without function
symbols. For \(n\ge1\), \(M_n\) is the set of all \(\nu\)-structures
on \(A=\{0,\ldots,n-1\}\). Fix \(p_R\in(0,1)\) for each relation
symbol. The measure \(\mu_n\) on \(M_n\) makes every ground relation
instance an independent Bernoulli event with parameter \(p_R\).
All subsets of this finite sample space are measurable.
The unbiased case has every \(p_R=1/2\).

Inside one fixed structure, a tuple of length \(\ell\) instead has
the uniform counting measure
\(\lambda_A^\ell(B)=|B|/n^\ell\) on \(A^\ell\).
Coordinates are sampled with replacement: syntactically distinct
variables may receive equal values. This tuple measure, not \(\mu_n\),
interprets a probability quantifier. The number
\(\mu_n(\theta)\) for a sentence is the probability of the model event
\(\{\mathcal A:\mathcal A\models\theta\}\).
\end{hbdefinition}

\begin{hbdefinition}{def:p15-pq-language}{The actual finitary formation and satisfaction clauses}
\(L_{\omega P}\) retains first-order equality, relation atoms, classical
Boolean connectives and ordinary quantifiers. If
\(\varphi(\bar x,\bar y)\) has exactly these free variables, with
the lists disjoint and without syntactic repetitions, then
\((\exists^{\ge r}\bar y)\varphi(\bar x,\bar y)\) is a formula,
binding \(\bar y\), for every real \(0<r<1\).
The source excludes probability quantifiers at \(0\) and \(1\).
Renaming and substitution are capture-avoiding.
For \(\bar a\in A^{|\bar x|}\),
\[
 \mathcal A\models(\exists^{\ge r}\bar y)\varphi(\bar a,\bar y)
 \quad\Longleftrightarrow\quad
 \frac{|\{\bar b\in A^{|\bar y|}:
                 \mathcal A\models\varphi(\bar a,\bar b)\}|}
      {n^{|\bar y|}}\ge r.
\]
Ordinary \(\forall\bar y,\exists\bar y\) abbreviate iterated
single-variable quantifiers; the probability binder counts tuples jointly.
Infinitary conjunctions belong to the separately named \(L_{\infty P}\);
they are not part of the fixed-size finitary translation below.
\end{hbdefinition}

\begin{hbtheorem}{thm:p15-pq-validities}{proved=annex:p15.pq.validities}{Semantic monotonicity and its boundary cases}
For these nonempty finite models, probability binders are monotone
in their scope and antitone in the threshold. For \(r,s\in(0,1)\),
\[
 \forall\bar y\,\varphi\ \Longrightarrow\
 (\exists^{\ge r}\bar y)\varphi\ \Longrightarrow\
 \exists\bar y\,\varphi.
\]
If \(r+s>1\), then
\((\exists^{\ge r}\bar y)\varphi\) implies
\(\neg(\exists^{\ge s}\bar y)\neg\varphi\).
Neither exchanging nested probability quantifiers nor replacing
negation by the complementary non-strict threshold is generally valid.
\end{hbtheorem}

\begin{hbexample}{def:p15-pq-order}{Order and threshold equality matter}
On three elements \(a,b,c\), let a binary relation have rows
\(\{a,b\},\{a,c\},\varnothing\), respectively.
Two rows contain at least two of the three elements, but only one column
does. Therefore
\((\exists^{\ge2/3}x)(\exists^{\ge2/3}y)R(x,y)\) is true,
whereas reversing the binders is false. On four elements with a unary
relation of size two, both half-threshold statements for that relation
and its complement are true. They are not each other's negations.
\end{hbexample}

\begin{hbtheorem}{thm:p15-pq-translation}{proved=annex:p15.pq.translation}{An exact size-indexed first-order translation}
For every fixed \(n\ge1\), finitary formulas have a capture-avoiding
first-order translation \(\tau_n\) preserving satisfaction in every
\(n\)-element structure and every parameter assignment.
A probability binder with \(\ell\) variables and threshold \(r\)
becomes the assertion that at least
\(k=\lceil rn^\ell\rceil\) pairwise distinct tuples satisfy the
translated scope. The annex gives the actual witness formula and both
directions, including the empty-tuple case.
This is a size-indexed translation, not uniform elimination on all
finite sizes. It is effective for rational thresholds; arbitrary real
threshold symbols have no assumed effective encoding.
\end{hbtheorem}

\begin{hbtheorem}{thm:p15-pq-limit-failures}{proved=annex:p15.pq.limit-failures}{Critical thresholds destroy unrestricted asymptotic claims}
For unbiased \(\mu_n\), the unary sentence
\((\exists^{\ge1/2}x)R(x)\) has limiting model probability \(1/2\),
not zero or one. With an unrestricted binary relation the sentence
\[
 \eta=\exists x\bigl[
  (\exists^{\ge1/2}y)R(x,y)\land
  (\exists^{\ge1/2}y)\neg R(x,y)\bigr]
\]
has probability zero at odd sizes and tends to one at even sizes.
Thus unrestricted \(L_{\omega P}\) has neither a zero--one law nor
a convergence law for this model distribution. This is the source's
Proposition~3.1 counterexample, with an independent complete proof.
The source's printed linear-exponent rate is not adopted; the annex
supplies a valid square-root-exponent bound and checks that distinction.
\end{hbtheorem}

The noncritical fragments, critical-value construction and elimination
proofs are developed separately in
\cref{sec:p15-probability-elimination}. They do not follow merely from
these basic clauses or \(\tau_n\), and their common-class, parameter
and size-dependent hypotheses must be retained.
No sound-and-complete calculus for the full language is claimed here.

The measure interfaces in \cref{def:p15-measure,def:p15-pushforward-law}
apply to the two finite measures once their sample spaces are declared.
P9's likelihood terms measure propositional world-events
(\cref{def:p09-likelihood-logic}); PCTL measures path events in a
transition model, with scheduler conventions on MDPs
(\cref{def:p09-pctl}). Pfeifer's primitive conditional-event prices and
coherence (\cref{def:p15-rat44-events,def:p15-rat44-coherence}) are different
again. None of their completeness, complexity or Bayes claims supplies
the quantified language treated here. Conversely this finite-model account
does not settle the original chapter's measure-model, completeness,
conditional-expectation or open-problem allocations.
